\section{Analysis and Discussion}
\label{sec:discussion}

\subsection{When Evaluation Dimensions Disagree}

The Baseball arm demonstrates why erasure cannot be certified with a single
task. An embedding edit directly alters the representations of a small token
set, so it can strongly change the likelihood of short answer options containing
those tokens. In contrast, concept exclusion removes learning signal from tens
of thousands of complete passages, which held-out text NLL measures more
directly. Moreover, Baseball's small raw question contrast is largely
reproduced by the unrelated no-Rome twin, whereas the chunk
difference-in-differences supplies an internal control for run-level drift.
Thus the instruments differ not only in what they emphasize but also in whether
they establish a concept-specific reference effect.

Discrete accuracy loses still more information. A change is counted only when
the gold option crosses a distractor, even if its likelihood falls on nearly
every example. This was not merely a reporting inconvenience. EMBER's delta
objective assigned specificity 1.0 at every candidate strength for both
concepts because neighboring-domain accuracy did not decline. For Ancient Rome,
the objective therefore rose to the largest tested delta, 200, without
bracketing an optimum. The continuous held-out evaluation later revealed both
overshoot and collateral damage. Baseball's efficacy happened to plateau at
delta 10, allowing the tie-break to select a gentler edit, but its specificity
signal remained equally uninformative.

\subsection{Concept-Dependent Erasure}

Beyond the measurement issue above, the two concept arms reveal genuine
heterogeneity. Baseball exclusion has 2.1 times Ancient Rome's text-NLL effect.
A retrieval-based locality audit suggests that the interventions remove
different kinds of evidence: Ancient Rome questions are more enriched for
supporting chunks that were blacklisted, whereas many Baseball questions ask
about rules repeated widely across the corpus. The removed Baseball slice also
contains highly predictable team and season prose. This helps explain why
chunk-NLL effects differ across concepts, but it is an interpretation rather
than a causal result.

EMBER's learned features also differ in quality. The selected Baseball feature
has a clearer concept-to-neutral activation ratio, and judged Baseball features
appear across far more factorization cells than Ancient Rome features. Coupled
with the much smaller selected delta, this may explain why Baseball avoids
Rome's heavy tail and general-text damage. It does not explain why the Baseball
update is almost orthogonal to the never-trained direction; localized editing
and distributed training remain fundamentally different interventions.

\subsection{Why Erasure Differs from Never-Training}

Concept exclusion changes every subsequent optimization step that would
otherwise have included a selected chunk. Its consequences can propagate
through embeddings, attention, MLPs, and later gradients on related text.
EMBER instead modifies 72--76 rows of one matrix after training. It is therefore
plausible for the two models to display a similar average target effect without
occupying nearby parameter states. Baseball makes this distinction explicit:
observable forgetting coexists with essentially zero alignment to the excluded
displacement.

This does not make parameter equality a necessary condition for behavioral
success; neural networks can implement similar functions with different
weights. It does make parameter-space comparison an informative mechanistic
test of the strong claim that erasure ``undoes'' learning. Under that reading,
EMBER does not retrace the training consequence in the completed arm.

RMU and SNMF will test whether a broader MLP intervention aligns more closely
with concept exclusion, or whether the mismatch is common to post-hoc editing.
\placeholder{Replace this prospective sentence with the completed comparison
and distinguish shared from method-specific failure modes.}
